\documentclass{amsart}
% For dealing with references we use the comment environment
\usepackage{verbatim}
\newenvironment{reference}{\comment}{\endcomment}
%\newenvironment{reference}{}{}
\newenvironment{slogan}{\comment}{\endcomment}
\newenvironment{history}{\comment}{\endcomment}

% For commutative diagrams we use Xy-pic
\usepackage[all]{xy}

% We use 2cell for 2-commutative diagrams.
\xyoption{2cell}
\UseAllTwocells

% We use multicol for the list of chapters between chapters
\usepackage{multicol}

% This is generally recommended for better output
\usepackage{lmodern}
\usepackage[T1]{fontenc}

% For cross-file-references

% Package for hypertext links:
\usepackage{hyperref}

% For any local file, say "hello.tex" you want to link to please

% Theorem environments.
%
\theoremstyle{plain}
\newtheorem{theorem}[subsection]{Theorem}
\newtheorem{proposition}[subsection]{Proposition}
\newtheorem{lemma}[subsection]{Lemma}

\theoremstyle{definition}
\newtheorem{definition}[subsection]{Definition}
\newtheorem{example}[subsection]{Example}
\newtheorem{exercise}[subsection]{Exercise}
\newtheorem{situation}[subsection]{Situation}

\theoremstyle{remark}
\newtheorem{remark}[subsection]{Remark}
\newtheorem{remarks}[subsection]{Remarks}

\numberwithin{equation}{subsection}

% Macros
%
\def\lim{\mathop{\mathrm{lim}}\nolimits}
\def\colim{\mathop{\mathrm{colim}}\nolimits}
\def\Spec{\mathop{\mathrm{Spec}}}
\def\Hom{\mathop{\mathrm{Hom}}\nolimits}
\def\Ext{\mathop{\mathrm{Ext}}\nolimits}
\def\SheafHom{\mathop{\mathcal{H}\!\mathit{om}}\nolimits}
\def\SheafExt{\mathop{\mathcal{E}\!\mathit{xt}}\nolimits}
\def\Sch{\mathit{Sch}}
\def\Mor{\mathop{\mathrm{Mor}}\nolimits}
\def\Ob{\mathop{\mathrm{Ob}}\nolimits}
\def\Sh{\mathop{\mathit{Sh}}\nolimits}
\def\NL{\mathop{N\!L}\nolimits}
\def\CH{\mathop{\mathrm{CH}}\nolimits}
\def\proetale{{pro\text{-}\acute{e}tale}}
\def\etale{{\acute{e}tale}}
\def\QCoh{\mathit{QCoh}}
\def\Ker{\mathop{\mathrm{Ker}}}
\def\Im{\mathop{\mathrm{Im}}}
\def\Coker{\mathop{\mathrm{Coker}}}
\def\Coim{\mathop{\mathrm{Coim}}}

% Boxtimes
%
\DeclareMathSymbol{\boxtimes}{\mathbin}{AMSa}{"02}

%
% Macros for moduli stacks/spaces
%
\def\QCohstack{\mathcal{QC}\!\mathit{oh}}
\def\Cohstack{\mathcal{C}\!\mathit{oh}}
\def\Spacesstack{\mathcal{S}\!\mathit{paces}}
\def\Quotfunctor{\mathrm{Quot}}
\def\Hilbfunctor{\mathrm{Hilb}}
\def\Curvesstack{\mathcal{C}\!\mathit{urves}}
\def\Polarizedstack{\mathcal{P}\!\mathit{olarized}}
\def\Complexesstack{\mathcal{C}\!\mathit{omplexes}}
% \Pic is the operator that assigns to X its picard group, usage \Pic(X)
% \Picardstack_{X/B} denotes the Picard stack of X over B
% \Picardfunctor_{X/B} denotes the Picard functor of X over B
\def\Pic{\mathop{\mathrm{Pic}}\nolimits}
\def\Picardstack{\mathcal{P}\!\mathit{ic}}
\def\Picardfunctor{\mathrm{Pic}}
\def\Deformationcategory{\mathcal{D}\!\mathit{ef}}

\setlength{\emergencystretch}{3em}
\begin{document}
\title[Stacks Project, Tag 0D4F]{463 --- Uniform projective reconstruction for flat proper finite-presentation thickenings with an ample lift}
\maketitle
\noindent Copyright (C) 2005--2025 Johan de Jong. Stacks Project authors. Source: \href{https://stacks.math.columbia.edu/tag/0D4F}{Tag 0D4F}.

\noindent Editorial difficulty note (not author text). Difficulty H2: One twist bound must work for every thickening in the statement, not just for a chosen deformation. The proof uses perfect cohomology and base change to transfer vanishing and local freeness, then proves finite presentation of the reconstructed graded algebra by Noetherian approximation while correcting the mismatch between two independent degree bounds..

\noindent Editorial provenance note (not author text). History: excerpt from revision \texttt{a0444\allowbreak 6e57e\allowbreak c1fbc\allowbreak 252a8\allowbreak 71afc\allowbreak ec775\allowbreak 2fb28\allowbreak 07b14}, more-morphisms.tex, lines 2852--3041. Reference presentation was adapted; original-author context and core support blocks were retained or added. The mathematical wording is unchanged. Licensed under GNU FDL 1.2 or later; no invariant sections or cover texts. See ../licenses/COPYING. Context blocks reproduce author definitions and supporting results; they are not additional counted problems.

\begin{definition}
\label{definition-thickening}
Thickenings.
\begin{enumerate}
\item We say a scheme $X'$ is a {\it thickening} of a scheme $X$ if
$X$ is a closed subscheme of $X'$ and the underlying topological spaces
are equal.
\item We say a scheme $X'$ is a {\it first order thickening} of a scheme $X$ if
$X$ is a closed subscheme of $X'$ and the quasi-coherent sheaf of ideals
$\mathcal{I} \subset \mathcal{O}_{X'}$ defining $X$ has square zero.
\item We say a scheme $X'$ is a {\it finite order thickening} of a scheme $X$
if $X$ is a closed subscheme of $X'$ and the quasi-coherent sheaf of ideals
$\mathcal{I} \subset \mathcal{O}_{X'}$ defining $X$ is nilpotent, i.e.,
there exists an integer $n \geq 0$ such that $\mathcal{I}^{n + 1} = 0$.
\item We say a scheme $X'$ is an {\it $n$th order thickening} of a scheme $X$
if $X$ is a closed subscheme of $X$ and $\mathcal{I}^{n + 1} = 0$
where $\mathcal{I} \subset \mathcal{O}_{X'}$ is the
quasi-coherent sheaf of ideals defining $X$.
\item Given two thickenings $X \subset X'$ and $Y \subset Y'$ a
{\it morphism of thickenings} is a morphism $f' : X' \to Y'$ such that
$f'(X) \subset Y$, i.e., such that $f'|_X$ factors through the closed
subscheme $Y$. In this situation we set $f = f'|_X : X \to Y$ and we say
that $(f, f') : (X \subset X') \to (Y \subset Y')$ is a morphism of
thickenings.
\item Let $S$ be a scheme. We similarly define {\it thickenings over $S$}, and
{\it morphisms of thickenings over $S$}. This means that the schemes
$X, X', Y, Y'$ above are schemes over $S$, and that the morphisms
$X \to X'$, $Y \to Y'$ and $f' : X' \to Y'$ are morphisms over $S$.
\end{enumerate}
\end{definition}

\begin{lemma}[Deformations of projective schemes]
\label{lemma-deform-projective}
Let $f : X \to S$ be a morphism of schemes which is proper, flat, and
of finite presentation. Let $\mathcal{L}$ be $f$-ample. Assume
$S$ is quasi-compact. There exists a $d_0 \geq 0$ such that
for every cartesian diagram
$$
\vcenter{
\xymatrix{
X \ar[r]_{i'} \ar[d]_f & X' \ar[d]^{f'} \\
S \ar[r]^i & S'
}
}
\quad\text{and}\quad
\begin{matrix}
\text{invertible }\mathcal{O}_{X'}\text{-module}\\
\mathcal{L}'\text{ with }\mathcal{L} \cong (i')^*\mathcal{L}'
\end{matrix}
$$
where $S \subset S'$ is a thickening and $f'$ is
proper, flat, of finite presentation we have
\begin{enumerate}
\item $R^p(f')_*(\mathcal{L}')^{\otimes d} = 0$
for all $p > 0$ and $d \geq d_0$,
\item $\mathcal{A}'_d = (f')_*(\mathcal{L}')^{\otimes d}$
is finite locally free for $d \geq d_0$,
\item $\mathcal{A}' =
\mathcal{O}_{S'} \oplus \bigoplus_{d \geq d_0} \mathcal{A}'_d$
is a quasi-coherent $\mathcal{O}_{S'}$-algebra of finite presentation,
\item there is a canonical isomorphism
$r' : X' \to \underline{\text{Proj}}_{S'}(\mathcal{A}')$, and
\item there is a canonical isomorphism
$\theta' : (r')^*\mathcal{O}_{\underline{\text{Proj}}_{S'}(\mathcal{A}')}(1)
\to \mathcal{L}'$.
\end{enumerate}
The construction of $\mathcal{A}'$, $r'$, $\theta'$
is functorial in the data $(X', S', i, i', f', \mathcal{L}')$.
\end{lemma}

\begin{proof}
We first describe the maps $r'$ and $\theta'$.
Observe that $\mathcal{L}'$ is $f'$-ample, see
Lemma \href{https://stacks.math.columbia.edu/tag/0D2R}{lemma thicken property relatively ample (Tag 0D2R)}.
There is a canonical map of quasi-coherent graded
$\mathcal{O}_{S'}$-algebras
$\mathcal{A}' \to \bigoplus_{d \geq 0} (f')_*(\mathcal{L}')^{\otimes d}$
which is an isomorphism in degrees $\geq d_0$.
Hence this induces an isomorphism on relative Proj
compatible with the Serre twists of the structure sheaf, see
Constructions, Lemma
\href{https://stacks.math.columbia.edu/tag/07ZJ}{lemma eventual iso graded rings map relative proj (Tag 07ZJ)}.
Hence we get the morphism $r'$ by
Morphisms, Lemma \href{https://stacks.math.columbia.edu/tag/01VJ}{lemma characterize relatively ample (Tag 01VJ)}
(which in turn appeals to the construction given in
Constructions, Lemma
\href{https://stacks.math.columbia.edu/tag/01O9}{lemma invertible map into relative proj (Tag 01O9)})
and it is an isomorphism by
Morphisms, Lemma \href{https://stacks.math.columbia.edu/tag/0C6J}{lemma proper ample is proj (Tag 0C6J)}.
We get the map $\theta'$ from Constructions, Lemma
\href{https://stacks.math.columbia.edu/tag/01O9}{lemma invertible map into relative proj (Tag 01O9)}.
By Properties, Lemma \href{https://stacks.math.columbia.edu/tag/01QI}{lemma ample gcd is one (Tag 01QI)}
we find that $\theta'$ is an isomorphism
(this also uses that the morphism $r'$ over affine
opens of $S'$ is the same as the morphism from
Properties, Lemma \href{https://stacks.math.columbia.edu/tag/01PZ}{lemma map into proj (Tag 01PZ)}
as is explained in the proof
of Morphisms, Lemma \href{https://stacks.math.columbia.edu/tag/0C6J}{lemma proper ample is proj (Tag 0C6J)}).

\medskip\noindent
Assuming the vanishing and local freeness stated in parts
(1) and (2), the functoriality of the construction can be seen as follows.
Suppose that $h : T \to S'$ is a morphism of schemes, denote
$f_T : X'_T \to T$ the base change of $f'$ and
$\mathcal{L}_T$ the pullback of $\mathcal{L}$ to $X'_T$.
By cohomology and base change
(as formulated in Derived Categories of Schemes,
Lemma \href{https://stacks.math.columbia.edu/tag/08IB}{lemma compare base change (Tag 08IB)} for example)
we have the corresponding vanishing over $T$ and moreover
$h^*\mathcal{A}'_d = f_{T, *}\mathcal{L}_T^{\otimes d}$
(and thus the local freeness of pushforwards as well
as the finite generation of the corresponding graded
$\mathcal{O}_T$-algebra $\mathcal{A}_T$).
Hence the morphism
$r_T : X_T \to
\underline{\text{Proj}}_T(\bigoplus f_{T, *}\mathcal{L}_T^{\otimes d})$
is simply the base change of $r'$ to $T$ and the pullback of
$\theta'$ is the map $\theta_T$.

\medskip\noindent
Having said all of the above, we see that it suffices to prove
(1), (2), and (3). Pick $d_0$ such that
$R^pf_*\mathcal{L}^{\otimes d} = 0$ for all $d \geq d_0$ and $p > 0$, see
Cohomology of Schemes, Lemma \href{https://stacks.math.columbia.edu/tag/0B5T}{lemma coherent proper ample (Tag 0B5T)}.
We claim that $d_0$ works.

\medskip\noindent
By cohomology and base change
(Derived Categories of Schemes,
Lemma \href{https://stacks.math.columbia.edu/tag/0B91}{lemma flat proper perfect direct image general (Tag 0B91)})
we see that $E'_d = Rf'_*(\mathcal{L}')^{\otimes d}$
is a perfect object of $D(\mathcal{O}_{S'})$
and its formation commutes with arbitrary base change.
In particular, $E_d = Li^*E'_d = Rf_*\mathcal{L}^{\otimes d}$.
By Derived Categories of Schemes, Lemma
\href{https://stacks.math.columbia.edu/tag/0D4E}{lemma vanishing implies locally free (Tag 0D4E)}
we see that for $d \geq d_0$ the complex $E_d$ is isomorphic to
the finite locally free $\mathcal{O}_S$-module
$f_*\mathcal{L}^{\otimes d}$ placed in
cohomological degree $0$. Then by
Derived Categories of Schemes, Lemma
\href{https://stacks.math.columbia.edu/tag/0BDK}{lemma open where cohomology in degree i rank r (Tag 0BDK)}
we conclude that $E'_d$ is isomorphic to a finite locally free
module placed in cohomological degree $0$.
Of course this means that $E'_d = \mathcal{A}'_d[0]$,
that $R^pf'_*(\mathcal{L}')^{\otimes d} = 0$ for $p > 0$,
and that $\mathcal{A}'_d$ is finite locally free.
This proves (1) and (2).

\medskip\noindent
The last thing we have to show is finite presentation of
$\mathcal{A}'$ as a sheaf of $\mathcal{O}_{S'}$-algebras
(this notion was introduced in Properties, Section
\href{https://stacks.math.columbia.edu/tag/01PD}{section extending quasi coherent sheaves (Tag 01PD)}).
Let $U' = \Spec(R') \subset S'$ be an affine open.
Then $A' = \mathcal{A}'(U')$ is a graded $R'$-algebra
whose graded parts are finite projective $R'$-modules.
We have to show that $A'$ is a finitely presented $R'$-algebra.
We will prove this by reduction to the Noetherian case.
Namely, we can find a finite type $\mathbf{Z}$-subalgebra
$R'_0 \subset R'$ and a pair\footnote{With the same properties
as those enjoyed by $X' \to S'$ and $\mathcal{L}'$, i.e.,
$X'_0 \to \Spec(R'_0)$ is flat and proper and $\mathcal{L}'_0$
is ample.} $(X'_0, \mathcal{L}'_0)$ over $R'_0$
whose base change is $(X'_{U'}, \mathcal{L}'|_{X'_{U'}})$, see
Limits, Lemmas
\href{https://stacks.math.columbia.edu/tag/01ZR}{lemma descend modules finite presentation (Tag 01ZR)},
\href{https://stacks.math.columbia.edu/tag/0B8W}{lemma descend invertible modules (Tag 0B8W)},
\href{https://stacks.math.columbia.edu/tag/081F}{lemma eventually proper (Tag 081F)},
\href{https://stacks.math.columbia.edu/tag/04AI}{lemma descend flat finite presentation (Tag 04AI)}, and
\href{https://stacks.math.columbia.edu/tag/09MT}{lemma limit ample (Tag 09MT)}.
Cohomology of Schemes, Lemma \href{https://stacks.math.columbia.edu/tag/0B5T}{lemma coherent proper ample (Tag 0B5T)}
implies
$A'_0 = \bigoplus_{d \geq 0} H^0(X'_0, (\mathcal{L}'_0)^{\otimes d})$
is a finitely generated graded $R'_0$-algebra and implies
there exists a $d'_0$ such that
$H^p(X'_0, (\mathcal{L}'_0)^{\otimes d}) = 0$, $p > 0$ for $d \geq d'_0$.
By the arguments given above applied to $X'_0 \to \Spec(R'_0)$ and
$\mathcal{L}'_0$ we see that $(A'_0)_d$ is a finite projective $R'_0$-module
and that
$$
A'_d = \mathcal{A}'_d(U') =
H^0(X'_{U'}, (\mathcal{L}')^{\otimes d}|_{X'_{U'}}) =
H^0(X'_0, (\mathcal{L}'_0)^{\otimes d}) \otimes_{R'_0} R' =
(A'_0)_d \otimes_{R'_0} R'
$$
for $d \geq d'_0$. Now a small twist in the argument is that we
don't know that we can choose $d'_0$ equal to $d_0$\footnote{Actually,
one can reduce to this case by doing more limit arguments.}. To
get around this we use the following sequence of arguments to finish
the proof:
\begin{enumerate}
\item[(a)] The algebra
$B = R'_0 \oplus \bigoplus_{d \geq \max(d_0, d'_0)} (A'_0)_d$
is an $R'_0$-algebra of finite type: apply
the Artin-Tate lemma to $B \subset A'_0$, see
Algebra, Lemma \href{https://stacks.math.columbia.edu/tag/00IS}{lemma Artin Tate (Tag 00IS)}.
\item[(b)] As $R'_0$ is Noetherian we see that
$B$ is an $R'_0$-algebra of finite presentation.
\item[(c)] By right exactness of tensor product we see that
$B \otimes_{R'_0} R'$ is an $R'$-algebra of finite presentation.
\item[(d)] By the displayed equalities this exactly says that
$C = R' \oplus \bigoplus_{d \geq \max(d_0, d'_0)} A'_d$
is an $R'$-algebra of finite presentation.
\item[(e)] The quotient $A'/C$ is the direct sum of the finite
projective $R'$-modules $A'_d$, $d_0 \leq d \leq \max(d_0, d'_0)$,
hence finitely presented as $R'$-module.
\item[(f)] The quotient $A'/C$ is finitely presented
as a $C$-module by Algebra, Lemma
\href{https://stacks.math.columbia.edu/tag/0561}{lemma finitely presented over subring (Tag 0561)}.
\item[(g)] Thus $A'$ is finitely presented as a $C$-module by
Algebra, Lemma \href{https://stacks.math.columbia.edu/tag/0519}{lemma extension (Tag 0519)}.
\item[(h)] By Algebra, Lemma \href{https://stacks.math.columbia.edu/tag/0D46}{lemma finite finite type (Tag 0D46)}
this implies $A'$ is finitely presented as a $C$-algebra.
\item[(i)] Finally, by
Algebra, Lemma \href{https://stacks.math.columbia.edu/tag/00F4}{lemma compose finite type (Tag 00F4)}
applied to $R' \to C \to A'$
this implies $A'$ is finitely presented as an $R'$-algebra.
\end{enumerate}
This finishes the proof.
\end{proof}
\end{document}
